%% 5-2 時間をどんどん短くしてみよう（第11回：ページ制限なし。偶数ページ）
\bkchapter{5}{2}{時間をどんどん短くしてみよう}{\bklevel{2}}{時間をどんどん短くしてみよう}
\bkmichishirube{3-8,3-9,5-1}{5-3,5-4}

\begin{goal}
\begin{itemize}[label={\color{bkTheme}●}, leftmargin=1.6\zw, itemsep=0pt, topsep=1pt]
\item \bklevel{2}\enspace $\Delta t$ を 0.1、0.01 と小さくして、平均の速さを表で計算できる。
\item \bklevel{2}\enspace 平均の速さを $h$ の式で表し、$h$ が小さいときの値を読める。
\item \bklevel{2}\enspace $x=at^2$ の形の運動で、瞬間の速さを求められる。
\end{itemize}
\end{goal}

\begin{hatena}
\noindent\begin{minipage}[c]{0.71\linewidth}
手からボールを落とすと、だんだん速くなりながら落ちていきます。落ちはじめてから $t$ 秒のあいだに落ちる距離は、
\[ x=5t^2\quad\text{（m）} \]
と表せます（ほんとうは $4.9t^2$ ですが、この章では計算しやすいように5にします）。1秒たった瞬間、ボールの速さはいくつでしょう。1秒で\qty{5}{m}落ちたから、秒速\qty{5}{m}？
\end{minipage}\hfill
\begin{minipage}[c]{0.26\linewidth}\raggedleft
\begin{tikzpicture}[x=1mm, y=0.9mm]
  \draw[bkGray, line width=0.8pt] (0,0) -- (0,-47);
  \foreach \t/\x in {0/0,1/5,2/20,3/45} {
    \draw[bkGray, line width=0.6pt] (-1.5,-\x) -- (1.5,-\x);
    \node[anchor=east, font=\sffamily\scriptsize] at (-2,-\x) {\x\ m};
    \shade[ball color=bkOrange] (8,-\x) circle (2.2);
    \node[anchor=west, font=\sffamily\scriptsize] at (11,-\x) {\t 秒};}
\end{tikzpicture}
\end{minipage}
\end{hatena}

図の数を、式で確かめてみましょう。2秒なら $5\times2^2=5\times4=20$、3秒なら $5\times3^2=5\times9=45$ です。1秒ごとに落ちる距離は、5、15、25 m と、どんどん大きくなっています。

5-1で調べたように、0秒から1秒までの平均の速さは $\dfrac{5-0}{1-0}=\qty{5}{m/s}$ です。でも、ボールはだんだん速くなっているので、1秒の瞬間の速さは、それより大きいはずです。この章では、あいだを短くして、瞬間の速さを追いかけます。

%% ---------------- やってみよう・思い出そう ----------------
\begin{yattemiyou}[表計算であいだを短くしよう]
表計算のシートに、次のように入力します。A列にあいだ $\Delta t$、B列に $1+\Delta t$ 秒の位置、C列に平均の速さを計算させます。
\begin{bkcenter}
\bksheet{3}{28}{%
1 & $\Delta t$ & $1+\Delta t$ 秒の位置 & 平均の速さ \\
2 & 1 & \texttt{=5*(1+A2)\^{}2} & \texttt{=(B2-5)/A2} \\
3 & 0.1 & （上をコピー） & （上をコピー） \\}
\end{bkcenter}
\begin{enumerate}[label={\bkcircled[bkGreen]{\arabic*}}, leftmargin=2.4\zw, itemsep=1pt, topsep=3pt]
\item A列に 1、0.1、0.01、0.001、0.0001 と入れ、B列とC列の式を下へコピーしよう。
\item C列の数は、何に近づいていくかな？
\end{enumerate}
\end{yattemiyou}

\noindent 式の意味を確かめておきます。\texttt{\^{}2} は2乗（2-18）です。B列の \texttt{=5*(1+A2)\^{}2} は、$x=5t^2$ の $t$ に $1+\Delta t$ を入れた「$1+\Delta t$ 秒の位置」です。C列の \texttt{=(B2-5)/A2} は、「あとの位置${}-{}$1秒の位置（5）」で進んだ距離を出し、それをあいだ $\Delta t$ でわった、平均の速さです。

\begin{bkcenter}
\bksheet{3}{28}{%
1 & $\Delta t$ & $1+\Delta t$ 秒の位置 & 平均の速さ \\
2 & 1 & 20 & 15 \\
3 & 0.1 & 6.05 & 10.5 \\
4 & 0.01 & 5.1005 & 10.05 \\
5 & 0.001 & 5.010005 & 10.005 \\
6 & 0.0001 & 5.00100005 & 10.0005 \\}
\bktabcaption{計算の結果。平均の速さは 10 に近づいていく。}
\end{bkcenter}

たとえば3行めは、$\Delta t=0.1$ なので1.1秒の位置 $5\times1.1^2=5\times1.21=6.05$ を求め、$(6.05-5)\div0.1=1.05\div0.1=10.5$ です。電卓でも確かめてみよう。

\begin{omoidasou}
\bkrecall{5-1}{平均の速さ $=\dfrac{\Delta x}{\Delta t}$。あいだを短くすると、瞬間の速さに近づく。$x(1.1)$ は「$t=1.1$ のときの $x$」。}
\bkrecall{3-9}{$y=ax^2$ のグラフは放物線。$x$ が2倍になると $y$ は4倍になる。}
\bkrecall{3-8}{$(a+b)^2=a^2+2ab+b^2$、$(a-b)^2=a^2-2ab+b^2$\\たとえば $(1+h)^2=1+2h+h^2$}
\bkrecall{2-12}{まとめてわるときは、1つずつわってよい：$(A+B)\div h=A\div h+B\div h$}
\end{omoidasou}

%% ---------------- 数学の道具 ----------------
\bksection{doguu}{数学の道具}

\bksubsection{1}{前からも後ろからも、10に近づく}

\leavevmode\label{bk:5-2:hyou}やってみようでは、1秒より後ろの時刻（$1+\Delta t$ 秒）をとりました。1秒より前の時刻（$1-\Delta t$ 秒）から1秒までの平均の速さも、計算してみましょう。たとえば $\Delta t=0.1$ なら、0.9秒の位置は $5\times0.9^2=5\times0.81=4.05$ です。0.9秒から1秒までに $5-4.05=0.95$（m）進んだので、
\[ \frac{5-4.05}{0.1}=\frac{0.95}{0.1}=9.5 \]

\begin{bkcenter}
\begin{bktable}{colspec={Q[l,m,wd=44mm] *{4}{Q[c,m,wd=16mm]}}, rowsep=1.5pt}
$\Delta t$ & 1 & 0.1 & 0.01 & 0.001 \\
1秒から $1+\Delta t$ 秒まで & 15 & 10.5 & 10.05 & 10.005 \\
$1-\Delta t$ 秒から1秒まで & 5 & 9.5 & 9.95 & 9.995 \\
\end{bktable}
\bktabcaption{1秒の前と後ろの、平均の速さ（m/s）}
\end{bkcenter}

\begin{bkfigure}
\begin{tikzpicture}[x=1mm, y=1mm]
  % 数直線（9.4〜10.6、1＝100 mm）。後ろからの値は上、前からの値は下に書き、高さをずらす
  \draw[bk arrow, bkInk, line width=0.8pt] (-2,0) -- (124,0);
  \foreach \v in {9.5,10,10.5} {\draw[bkInk, line width=0.6pt] ({(\v-9.4)*100},-1.2) -- ({(\v-9.4)*100},1.2);}
  \fill[bkOrange] (60,0) circle (1.3);
  \node[anchor=south, font=\sffamily\bfseries\small, text=bkOrange!80!black] at (59,2.2) {10};
  \foreach \v/\hh in {10.5/4, 10.05/10} {\fill[bkBlue] ({(\v-9.4)*100},0) circle (0.9);
    \draw[bkBlue, line width=0.5pt] ({(\v-9.4)*100},0) -- ({(\v-9.4)*100},\hh-0.6);
    \node[anchor=south, font=\sffamily\scriptsize, text=bkBlue] at ({(\v-9.4)*100},\hh) {\v};}
  \foreach \v/\hh in {9.5/4, 9.95/10} {\fill[bkGreen!70!black] ({(\v-9.4)*100},0) circle (0.9);
    \draw[bkGreen!70!black, line width=0.5pt] ({(\v-9.4)*100},0) -- ({(\v-9.4)*100},-\hh+0.6);
    \node[anchor=north, font=\sffamily\scriptsize, text=bkGreen!70!black] at ({(\v-9.4)*100},-\hh) {\v};}
  \draw[bk arrow, bkBlue, line width=0.9pt] (118,16) -- (70,16);
  \node[anchor=south, font=\sffamily\footnotesize, text=bkBlue] at (94,17) {後ろから（10より大きい側）};
  \draw[bk arrow, bkGreen!70!black, line width=0.9pt] (2,-16) -- (50,-16);
  \node[anchor=north, font=\sffamily\footnotesize, text=bkGreen!70!black] at (26,-17) {前から（10より小さい側）};
\end{tikzpicture}
\bkfigcaption{数直線で見ると、両側から10をはさみながら近づいていく。}
\end{bkfigure}

後ろからは10より少し大きい値で、前からは10より少し小さい値で、どちらも10に近づいていきます。10を両側からはさみうちにしているので、近づく先は10しかありません。だから、\textbf{1秒の瞬間の速さは\qty{10}{m/s}}だと考えられます。

\bkhatenaref の「秒速\qty{5}{m}？」は、0秒から1秒までの平均の速さでした。瞬間の速さ10はその2倍です。なぜ2倍なのでしょう。あとでわかるように、このボールの速さは0秒で0、1秒で10と、時間に比例して同じ割合でふえていきます。同じ割合でふえるとき、0から10までをならした平均は、まん中の $(0+10)\div2=5$ です。だから、平均の5と瞬間の10は、ちょうど2倍のちがいになります。

\bksubsection{2}{文字 $h$ を使って、近づき方を式にする}

\leavevmode\label{bk:5-2:h}表では、$\Delta t$ を変えるたびに計算し直しました。0.1でも、0.01でも、同じ計算のくり返しです。そこで、あいだを文字 $h$ で表して、1回の計算ですませましょう。$h$ は、0.1 や 0.01 など、どんな数でも入れられる箱です（2-11）。ただし、あとで $h$ でわるので、$h$ は0でない数とします。

\noindent\textbf{①\enspace $1+h$ 秒の位置}\par
$x=5t^2$ の $t$ に、かっこごと $1+h$ を入れます（代入、2-13）。
\[ x(1+h)=5(1+h)^2 \]

\noindent\textbf{②\enspace かっこをはずす（展開）}\par
$(1+h)^2$ は、1辺が $1+h$ の正方形の面積です。図のように4つに分けると、
\[ (1+h)^2=1+h+h+h^2=1+2h+h^2 \]

\noindent\begin{minipage}[c]{0.52\linewidth}
\begin{tikzpicture}[x=1mm, y=1mm]
  % 1辺 1+h の正方形（1→34 mm、h→9 mm。h は見やすいように大きめにかいた）
  \fill[bkBlue!14] (0,9) rectangle (34,43);
  \fill[bkOrange!25] (34,9) rectangle (43,43);
  \fill[bkOrange!25] (0,0) rectangle (34,9);
  \fill[bkRed!35] (34,0) rectangle (43,9);
  \draw[bkInk, line width=0.8pt] (0,0) rectangle (43,43);
  \draw[bkInk, line width=0.6pt] (34,0) -- (34,43) (0,9) -- (43,9);
  \node[font=\sffamily\small] at (17,26) {$1\times1=1$};
  \node[font=\sffamily\small] at (38.5,26) {$h$};
  \node[font=\sffamily\small] at (17,4.5) {$h$};
  \node[font=\sffamily\scriptsize] at (38.5,4.5) {$h^2$};
  \node[anchor=south, font=\sffamily\footnotesize] at (17,44.5) {1};
  \node[anchor=south, font=\sffamily\footnotesize] at (38.5,44.5) {$h$};
  \node[anchor=east, font=\sffamily\footnotesize] at (-1.5,26) {1};
  \node[anchor=east, font=\sffamily\footnotesize] at (-1.5,4.5) {$h$};
\end{tikzpicture}
\end{minipage}\hfill
\begin{minipage}[c]{0.45\linewidth}\raggedright
大きな正方形 1 と、細長い長方形 $1\times h=h$ が2まい、小さな正方形 $h\times h=h^2$ が1まいです。

$h=0.1$ で確かめると、$1.1^2=1.21$、\\$1+2\times0.1+0.01=1.21$ で合います。
\end{minipage}

\medskip
だから、$x(1+h)=5(1+2h+h^2)$。かっこの中の1つずつに5をかけて（分配法則、2-12）、
\[ x(1+h)=5+10h+5h^2 \]

\noindent\textbf{③\enspace 進んだ距離 $\Delta x$（あと${}-{}$まえ）}\par
1秒の位置は $x(1)=5$ なので、
\[ \Delta x=(5+10h+5h^2)-5=10h+5h^2 \]

\noindent\textbf{④\enspace $h$ でわる}\par
まとめてわるときは、1つずつわってよいので（2-12）、
\[ \frac{\Delta x}{h}=\frac{10h+5h^2}{h}=\frac{10h}{h}+\frac{5h^2}{h}=10+5h \]
$\dfrac{10h}{h}$ は、分子と分母の $h$ を約分して10です。$\dfrac{5h^2}{h}=\dfrac{5\times h\times h}{h}$ は、$h$ を1つ約分して $5h$ です。

\medskip
こうして、1秒から $1+h$ 秒までの平均の速さは、
\[ \frac{\Delta x}{\Delta t}=10+5h \]
という式になりました。表と合っているか、確かめましょう。
\begin{align*}
h=0.1\ \text{なら}&\quad 10+5\times0.1=10+0.5=10.5\\
h=0.01\ \text{なら}&\quad 10+5\times0.01=10+0.05=10.05
\end{align*}
表の10.5、10.05とぴったり合います。

\begin{bkfigure}
\begin{tikzpicture}[x=1mm, y=1mm]
  % 第6回：左に行の見出し、箱の中は数だけ。矢印は箱の名前どうしでつなぎ、両端を離す
  \node[anchor=east, font=\sffamily\small, text=bkInk, inner sep=0pt] at (14,3.2) {$h$};
  \node[anchor=east, font=\sffamily\small, text=bkInk, inner sep=0pt] at (14,-3.2) {$10+5h$};
  \foreach \h/\v/\xx/\nm in {1/15/27/b1, 0.1/10.5/61/b2, 0.01/10.05/95/b3, 0.001/10.005/129/b4} {
    \node[bk node, minimum width=22mm, inner xsep=2mm, font=\sffamily\small] (\nm) at (\xx,0) {$\h$\\[1pt]$\v$};}
  \foreach \a/\b in {b1/b2,b2/b3,b3/b4} \draw[bk link] (\a) -- (\b);
  \draw[bk link, bkOrange, shorten <=1.2mm] (b4.south) -- ++(0,-6.5) coordinate (ar);
  \node[anchor=north, font=\sffamily\bfseries\small, text=bkOrange!80!black, inner sep=0pt] at ([yshift=-1.5mm]ar) {10 に近づく};
  \node[anchor=north west, font=\sffamily\small, text=bkInk, inner sep=0pt] at (16,-9.5) {$h$ が小さくなると、$5h$ の部分が、どんどん小さくなる};
\end{tikzpicture}
\end{bkfigure}

$h$ を小さくすると、$10+5h$ の $5h$ の部分がどんどん0に近づき、全体は10に近づきます。表では「10に近づいていそう」としか言えませんでしたが、式を見れば、近づく先が10だとはっきりわかります。

\noindent\textbf{10と $5h$ は、何を表している？}\par
$\Delta x=10h+5h^2$ の2つの部分にも、意味があります。$10h$ は、「1秒のときの速さ\qty{10}{m/s}のまま、$h$ 秒走ったときの距離」です（距離＝速さ×時間）。$5h^2$ は、そのあいだにボールが速くなった分だけ、よけいに進んだ距離です。$h$ が小さいと、$h^2$ は $h$ よりずっと小さくなります（$h=0.01$ なら $h^2=0.0001$）。だから、短いあいだでは「よけいな分」はほとんど消えて、平均の速さは10にほぼ等しくなるのです。

\bksubsection{3}{どの時刻でも計算できる}

\leavevmode\label{bk:5-2:t}1秒のかわりに、どの時刻 $t$ でも同じ計算ができます。$t$ 秒から $t+h$ 秒までを考えます。ここでは $t$ は「いま考えている時刻」で、計算のあいだは決まった数として扱います。$h$ は、そこからのあいだです。

②の正方形の図の「1」を「$t$」にかえると、$(t+h)^2=t^2+2th+h^2$ です（3-8）。①〜④と同じ段で、1行ずつ計算します（1・2行め：代入して展開、3行め：あと${}-{}$まえ、4行め：1つずつ $h$ でわる）。
\begin{align*}
x(t+h)&=5(t+h)^2=5(t^2+2th+h^2)\\
&=5t^2+10th+5h^2\\
\Delta x&=(5t^2+10th+5h^2)-5t^2=10th+5h^2\\
\frac{\Delta x}{h}&=\frac{10th+5h^2}{h}=10t+5h
\end{align*}
$t=1$ を入れると $10+5h$ で、さっきの式と同じになります。$h$ を0に近づけると $5h$ が消えて、$t$ 秒の瞬間の速さ $v$ は
\[ v=10t\quad\text{（m/s）} \]
です（$v$ は速さを表す文字。velocity の頭の文字、2-11）。

\begin{bkcenter}
\begin{bktable}{colspec={Q[l,m,wd=52mm] *{5}{Q[c,m,wd=9mm]}}, rowsep=1.5pt}
時刻 $t$（s） & 0 & 1 & 2 & 3 & 4 \\
位置 $x=5t^2$（m） & 0 & 5 & 20 & 45 & 80 \\
瞬間の速さ $v=10t$（m/s） & 0 & 10 & 20 & 30 & 40 \\
\end{bktable}
\bktabcaption{落ちるボールの位置と、瞬間の速さ}
\end{bkcenter}

瞬間の速さは、1秒ごとに\qty{10}{m/s}ずつ増えています。2-0で「落ちるボールの速さは、1秒ごとに約\qty{10}{m/s}ずつふえる」と書いたのは、このことです。式から、それが確かめられました。

\bksubsection{4}{$x=at^2$ なら $v=2at$}

\leavevmode\label{bk:5-2:a}5のかわりに、どんな数 $a$ でも同じ計算ができます。$x=at^2$ とします（$a$ は決まった数）。
\begin{align*}
x(t+h)-x(t)&=a(t+h)^2-at^2\\
&=at^2+2ath+ah^2-at^2=2ath+ah^2\\
\frac{\Delta x}{h}&=\frac{2ath+ah^2}{h}=2at+ah
\end{align*}
$h$ を0に近づけると $ah$ が消えて、
\[ v=2at \]
です。2という数は、展開の $(t+h)^2=t^2+\underline{2th}+h^2$ の「2」から来ています。正方形の図で、細長い長方形が\textbf{2まい}あったからです。

\begin{bkcenter}
\begin{bktable}{colspec={Q[c,m,wd=26mm] Q[c,m,wd=20mm] Q[c,m,wd=40mm]}, rowsep=1.5pt}
位置の式 & $a$ & 瞬間の速さ $v=2at$ \\
$x=5t^2$ & 5 & $v=2\times5\times t=10t$ \\
$x=3t^2$ & 3 & $v=2\times3\times t=6t$ \\
$x=20t^2$ & 20 & $v=2\times20\times t=40t$ \\
$x=4.9t^2$ & 4.9 & $v=2\times4.9\times t=9.8t$ \\
\end{bktable}
\bktabcaption{$x=at^2$ の瞬間の速さ}
\end{bkcenter}

3行めの $x=20t^2$ は、5-1のボールです。$v=40t$ に $t=1$ を入れると40で、5-1で表から見つけた\qty{40}{cm/s}と合います。

4行めの $x=4.9t^2$ が、ほんとうの落ちる物の式です。$v=9.8t$ なので、落ちる物は、1秒ごとに\qty{9.8}{m/s}ずつ速くなります（5-5）。\bkhatenaref で「ほんとうは4.9」と書いた4.9は、この9.8の半分です。

\noindent\textbf{同じ速さで動く物なら}\par
位置が $x=20t$（m）の物は、どうでしょう。
\[ \frac{x(t+h)-x(t)}{h}=\frac{20(t+h)-20t}{h}=\frac{20t+20h-20t}{h}=\frac{20h}{h}=20 \]
$h$ がどんな数でも、平均の速さは20のままです。グラフがまっすぐな直線なので、どのあいだをとっても傾きは20です。だから、瞬間の速さもいつも\qty{20}{m/s}です。

\bksubsection{5}{グラフで見ると}

\noindent\begin{minipage}[c]{0.5\linewidth}
\begin{tikzpicture}
\begin{axis}[bkplot, width=66mm, height=48mm, xmin=0, xmax=3.4, ymin=0, ymax=56, xtick={0,1,2,3}, ytick={0,10,...,50},
  xlabel={$t$（s）}, ylabel={$x$（m）}]
  % 第6回：接線は細くして曲線の上に重ね、曲線が見えるようにする。接線は枠（ymin〜ymax）の中で止める
  \addplot[bkcurve, domain=0:3.1] {5*x^2};
  \foreach \t/\a/\b in {1/0.56/2.0, 2/1.25/2.9, 3/2.1/3.36} {\edef\tmp{\noexpand\addplot[bkcurve2, line width=1pt, domain=\a:\b] {5*\t*\t+10*\t*(x-\t)};}\tmp}
  \addplot[only marks, mark=*, mark size=1.6pt, bkInk] coordinates {(1,5) (2,20) (3,45)};
\end{axis}
\end{tikzpicture}
\end{minipage}\hfill
\begin{minipage}[c]{0.46\linewidth}\raggedright
$x$-$t$ グラフの各点に、瞬間の速さを傾きにもつ直線（オレンジ）を引きました。傾きは、1秒で10、2秒で20、3秒で30。グラフがだんだん急になっていくようすが、傾きの数で表せました。
\end{minipage}

\bksubsection{6}{前からの近づき方も、式で確かめる}

$1-h$ 秒から1秒までの平均の速さも、$(a-b)^2=a^2-2ab+b^2$（3-8）を使って、式にしてみましょう。こんどは1秒が「あと」、$1-h$ 秒が「まえ」です。
\begin{align*}
x(1-h)&=5(1-h)^2=5(1-2h+h^2)=5-10h+5h^2 \\
\Delta x&=5-(5-10h+5h^2)=5-5+10h-5h^2=10h-5h^2 \\
\frac{\Delta x}{h}&=\frac{10h-5h^2}{h}=10-5h
\end{align*}
2行めでは、マイナスのかっこをはずすと、中の符号が全部かわることを使いました（2-12）。$h=0.1$ なら $10-0.5=9.5$ で、表と合います。

後ろからは $10+5h$、前からは $10-5h$。$h$ を0に近づけると、どちらも10になります。近づく先が同じなので、「1秒の瞬間の速さは10」と決められます。

\begin{taisetsu}
平均の速さを $h$ の式で表し、$h$ を0に近づける。残った値が瞬間の速さ。\\
$x=5t^2$ のとき、平均の速さは $10t+5h$、瞬間の速さは $v=10t$\\
$x=at^2$ のとき、瞬間の速さは $v=2at$
\end{taisetsu}

\begin{kotoba}
\small
\bkword{$h$}{あいだの長さを表す文字。0ではない、小さい数。}
\bkword{近づく}{$h$ を0に近づけたとき、式の値がある数にどんどん近くなること。5-3で「極限」という名前をつける。}
\bkword{$v$}{速さを表す文字（velocity の頭の文字）。}
\end{kotoba}

%% ---------------- 例題 ----------------
\bksection{reidai}{例題でたしかめよう}

\begin{reidai}[表で近づき方を見る]
$x=5t^2$（m）のボールについて、2秒から $2+\Delta t$ 秒までの平均の速さを、$\Delta t=0.1$ と $0.01$ で求め、2秒の瞬間の速さを答えましょう。\par
\kangaekata はじめに2秒の位置を求めます。次に、後ろの時刻の位置を計算して、（あと${}-{}$まえ）÷あいだ を求めます。\par
\shiki $x(2)=5\times2^2=5\times4=20$\par
\hspace*{2.6\zw}$x(2.1)=5\times2.1^2=5\times4.41=22.05$\par
\hspace*{2.6\zw}$x(2.01)=5\times2.01^2=5\times4.0401=20.2005$\par
\hspace*{2.6\zw}$(22.05-20)\div0.1=2.05\div0.1=20.5$\par
\hspace*{2.6\zw}$(20.2005-20)\div0.01=0.2005\div0.01=20.05$\par
\kotae 20.5、20.05と20に近づくので、\qty{20}{m/s}
\end{reidai}

\begin{jibun}
\bkreidairef{1}のボールの3秒の瞬間の速さを、同じ方法で求めましょう。
\bkjans{$x(3)=45$。$x(3.1)=48.05$ より 30.5、$x(3.01)=45.3005$ より 30.05。\qty{30}{m/s}}
\end{jibun}

\begin{reidai}[$h$ の式で表す]
$x=5t^2$ のボールについて、2秒から $2+h$ 秒までの平均の速さを $h$ の式で表し、2秒の瞬間の速さを求めましょう。\par
\kangaekata 代入して展開してから、「あと${}-{}$まえ」をとり、$h$ でわります。$(2+h)^2=4+4h+h^2$（正方形の図の1を2にかえる）。\par
\shiki $x(2+h)=5(2+h)^2=5(4+4h+h^2)=20+20h+5h^2$\par
\hspace*{2.6\zw}$\Delta x=(20+20h+5h^2)-20=20h+5h^2$\par
\hspace*{2.6\zw}$\dfrac{\Delta x}{h}=\dfrac{20h+5h^2}{h}=20+5h$\par
\kotae 平均の速さは $20+5h$。$h$ を0に近づけて、瞬間の速さは\qty{20}{m/s}
\end{reidai}

\begin{jibun}
$x=5t^2$ について、3秒から $3+h$ 秒までの平均の速さを $h$ の式で表し、3秒の瞬間の速さを求めましょう。
\bkjans{$5(3+h)^2=45+30h+5h^2$ より $30+5h$。\qty{30}{m/s}}
\end{jibun}

\begin{reidai}[ちがう式の運動]
位置が $x=3t^2$（m）で表される物体の、1秒の瞬間の速さを求めましょう。\par
\kangaekata 5のかわりに3になっただけです。同じ4つの段（代入・展開、あと${}-{}$まえ、$h$ でわる、$h$ を0に近づける）で計算します。\par
\shiki $x(1)=3\times1^2=3$\par
\hspace*{2.6\zw}$x(1+h)=3(1+h)^2=3(1+2h+h^2)=3+6h+3h^2$\par
\hspace*{2.6\zw}$\Delta x=(3+6h+3h^2)-3=6h+3h^2$\par
\hspace*{2.6\zw}$\dfrac{\Delta x}{h}=\dfrac{6h+3h^2}{h}=6+3h$\par
\kotae $h$ を0に近づけて、\qty{6}{m/s}（$v=2at=2\times3\times1=6$ とも合う）
\end{reidai}

\begin{jibun}
位置が $x=2t^2$（m）の物体の、3秒の瞬間の速さを求めましょう。
\bkjans{$2(3+h)^2=18+12h+2h^2$ より $12+2h$。\qty{12}{m/s}}
\end{jibun}

\begin{reidai}[項が2つある式]
位置が $x=2t^2+t$（m）の物体の、1秒の瞬間の速さを求めましょう。\par
\kangaekata $t$ が2か所にあります。どちらの $t$ にも、かっこごと $1+h$ を入れることを忘れずに。\par
\shiki $x(1)=2\times1^2+1=3$\par
\hspace*{2.6\zw}$x(1+h)=2(1+h)^2+(1+h)$\par
\hspace*{2.6\zw}$\phantom{x(1+h)}=2+4h+2h^2+1+h=3+5h+2h^2$\par
\hspace*{2.6\zw}$\Delta x=(3+5h+2h^2)-3=5h+2h^2$\par
\hspace*{2.6\zw}$\dfrac{\Delta x}{h}=\dfrac{5h+2h^2}{h}=5+2h$\par
\kotae $h$ を0に近づけて、\qty{5}{m/s}
\end{reidai}

\begin{jibun}
位置が $x=t^2+3t$（m）の物体の、2秒の瞬間の速さを求めましょう。
\bkjans{$x(2)=10$、$x(2+h)=10+7h+h^2$ より $7+h$。\qty{7}{m/s}}
\end{jibun}

\begin{reidai}[$t$ 秒の速さの式を使う]
$x=5t^2$（m）のボールの $t$ 秒の瞬間の速さを $t$ の式で表し、速さが\qty{30}{m/s}になるのは何秒後か求めましょう。そのとき、何m落ちていますか。\par
\kangaekata $t$ 秒から $t+h$ 秒までの平均の速さは $10t+5h$ でした（\pageref{bk:5-2:t}ページ）。$h$ を0に近づけて速さの式を作り、それが30になる $t$ を求めます。\par
\shiki $h$ を0に近づけて $v=10t$\par
\hspace*{2.6\zw}$10t=30$。両辺を10でわって（2-14）$t=3$\par
\hspace*{2.6\zw}$x=5\times3^2=5\times9=45$\par
\noindent\begin{minipage}[c]{0.5\linewidth}
\begin{tikzpicture}
\begin{axis}[bkplot, width=56mm, height=38mm, xmin=0, xmax=4.3, ymin=0, ymax=44, xtick={0,1,2,3,4}, ytick={0,10,20,30,40},
  xlabel={$t$（s）}, ylabel={$v$（m/s）}]
  \addplot[bkcurve2, domain=0:4] {10*x};
  \draw[bkGray, line width=0.6pt, dash pattern=on 2pt off 1.5pt] (axis cs:0,30) -- (axis cs:3,30) -- (axis cs:3,0);
  \addplot[only marks, mark=*, mark size=1.7pt, bkInk] coordinates {(3,30)};
\end{axis}
\end{tikzpicture}
\end{minipage}\hfill
\begin{minipage}[c]{0.46\linewidth}\raggedright
$v=10t$ のグラフ（速さと時間のグラフ）。たて軸の30から横に進み、グラフに当たった所の $t$ を読むと3秒です。
\end{minipage}\par
\kotae $v=10t$、3秒後、\qty{45}{m}
\end{reidai}

\begin{jibun}
$x=2t^2$（m）の物体の $t$ 秒の瞬間の速さを $t$ の式で表し、速さが\qty{12}{m/s}になる時刻と、そのときの位置を求めましょう。
\bkjans{$2(t+h)^2-2t^2=4th+2h^2$ より $v=4t$。$4t=12$ より3秒後、$x=\qty{18}{m}$}
\end{jibun}

\begin{reidai}[上に投げたボール（速さがマイナスになる）]
ボールを真上に投げました。手の位置を0、上向きを正として、$t$ 秒後の高さは $x=10t-5t^2$（m）です。1.5秒の瞬間の速さを求めましょう。\par
\kangaekata $-5t^2$ の項にも、かっこごと入れます。マイナスのかっこをはずすときは、中の符号が全部かわります（2-12）。\par
\shiki $x(1.5)=10\times1.5-5\times2.25=15-11.25=3.75$\par
\hspace*{2.6\zw}$x(1.5+h)=10(1.5+h)-5(1.5+h)^2$\par
\hspace*{2.6\zw}$\phantom{x(1.5+h)}=15+10h-5(2.25+3h+h^2)$\par
\hspace*{2.6\zw}$\phantom{x(1.5+h)}=15+10h-11.25-15h-5h^2=3.75-5h-5h^2$\par
\hspace*{2.6\zw}$\Delta x=(3.75-5h-5h^2)-3.75=-5h-5h^2$\par
\hspace*{2.6\zw}$\dfrac{\Delta x}{h}=\dfrac{-5h-5h^2}{h}=-5-5h$\par
\kotae $h$ を0に近づけて、$\qty{-5}{m/s}$。マイナスは下向き。1.5秒には、ボールは\qty{5}{m/s}で落ちてきている。
\end{reidai}

\begin{jibun}
\bkreidairef{6}のボールの、0.5秒の瞬間の速さを求めましょう。上向きと下向きのどちらに動いていますか。
\bkjans{$x(0.5)=3.75$、$x(0.5+h)=3.75+5h-5h^2$ より $5-5h$。\qty{5}{m/s}で上向き}
\end{jibun}

\begin{reidai}[$v=2at$ を使う]
5-1の坂をころがるボールの位置は $x=20t^2$（cm）でした。$v=2at$ を使って $t$ 秒の瞬間の速さを式で表し、速さが\qty{100}{cm/s}になる時刻と、そのときの位置を求めましょう。\par
\kangaekata $x=at^2$ の $a$ が20です。$v=2at$ の $a$ に20を入れます。\par
\shiki $v=2\times20\times t=40t$\par
\hspace*{2.6\zw}$40t=100$。両辺を40でわって $t=2.5$\par
\hspace*{2.6\zw}$x=20\times2.5^2=20\times6.25=125$\par
\kotae $v=40t$（cm/s）、2.5秒、\qty{125}{cm}
\end{reidai}

\begin{jibun}
ほんとうの落ちる物の式は $x=4.9t^2$（m）です。$v=2at$ を使って $t$ 秒の瞬間の速さを式で表し、速さが\qty{19.6}{m/s}になる時刻と、それまでに落ちた距離を求めましょう。
\bkjans{$v=9.8t$。$9.8t=19.6$ より2秒、$4.9\times2^2=\qty{19.6}{m}$}
\end{jibun}

\begin{column}{ガリレオは、どうやって時間をはかった？}
ガリレオのころ（400年ほど前）には、ストップウォッチはありませんでした。ガリレオは、大きな水おけの底に細い管をつけ、球が坂をころがっているあいだだけ水を流し、出てきた水の重さをはかりました。大きな水おけなので、水面の高さはほとんど変わらず、水はほぼ同じ速さで流れ出ます。だから、水の重さは時間に比例します。水の重さが2倍なら、時間も2倍です。こうして、「進む距離は時間の2乗に比例する」ことを確かめたのです。「はかれない量を、はかれる量に置きかえる」という、計測の工夫です。
\end{column}

\clearpage
%% ---------------- まちがえやすいところ・物理工学 ----------------
\bksection{tsumazuki}{まちがえやすいところをチェック}

\begin{tsumazuki}{いきなり $h=0$ を入れてしまう}
\vspace{7pt}% 1行目が分数で高いので、題とのあいだをほかの箱とそろえる
$\dfrac{10h+5h^2}{h}$ に $h=0$ を入れると $\dfrac{0}{0}$ になり、計算できません。0でわることはできないからです（1-14）。先に $h$ でわって $10+5h$ にしてから、$h$ を0に近づけます。$h$ そのものは、最後まで0ではありません。
\end{tsumazuki}

\begin{tsumazuki}{展開で $2h$ をわすれる}
\bkng{$(1+h)^2=1+h^2$}\bkok{$(1+h)^2=(1+h)(1+h)=1+h+h+h^2=1+2h+h^2$}
1辺 $1+h$ の正方形を4つに分けた図を思い出すと、まちがえません（3-8）。$h=0.1$ で確かめると、$1.1^2=1.21$ ですが、$1+0.1^2=1.01$ で合いません。
\end{tsumazuki}

\begin{tsumazuki}{「あと${}-{}$まえ」の「まえ」をひき忘れる}
\bkng{$\dfrac{x(1+h)}{h}=\dfrac{5+10h+5h^2}{h}$}\bkok{$\dfrac{x(1+h)-x(1)}{h}=\dfrac{10h+5h^2}{h}$}
平均の速さは、進んだ距離（変化量）÷時間です。$x(1+h)$ は「位置」で、進んだ距離ではありません。
\end{tsumazuki}

\begin{tsumazuki}{2か所の $t$ の片方にだけ入れる}
\bkng{$x=2t^2+t$ で、$x(1+h)=2(1+h)^2+1$}\bkok{$x(1+h)=2(1+h)^2+(1+h)$}
$t$ がある所は、全部かっこごと入れかえます（2-13）。
\end{tsumazuki}

\begin{butsuri}
\begin{itemize}[leftmargin=1.2\zw, itemsep=2pt, topsep=0pt, label={\color{bkPurple}●}]
\item \textbf{ガリレオの斜面の実験}\par
400年ほど前、ガリレオは坂をころがる球を調べ、進む距離が時間の2乗に比例する（$x=at^2$）ことを見つけました。この章の計算から、速さは $v=2at$ で、時間に比例して増えることがわかります。
\item \textbf{車の加速性能}\par
「時速\qty{100}{km}まで○秒」という性能は、瞬間の速さがどれだけ早く増えるかを表しています。
\item \textbf{落下の衝撃}\par
$v=9.8t$ なので、落ちる時間が長いほど速くなります。2秒落ちると約\qty{20}{m/s}（時速約\qty{70}{km}）にもなります。高い所から物を落とさないのは、このためです。
\end{itemize}
\end{butsuri}

\begin{column}{表計算で $h$ を小さくしすぎると}
やってみようの表計算で、$h$ をもっと小さくしてみましょう。B列は、やってみようと同じ「（あと${}-{}$まえ）÷あいだ」で計算した平均の速さ、C列は、$h$ の式 $10+5h$ で計算した値です。
\begin{bkcenter}
\bksheet{3}{30}{%
1 & $h$ & 平均の速さ & $10+5h$ \\
2 & 0.1 & 10.5 & 10.5 \\
3 & 0.001 & 10.005 & 10.005 \\
4 & 0.00001 & 10.00005 & 10.00005 \\
5 & $10^{-8}$ & {\color{bkRed}\bfseries 9.9999999392…} & 10.00000005 \\
6 & $10^{-10}$ & {\color{bkRed}\bfseries 10.0000008274…} & 10.0000000005 \\
7 & $10^{-12}$ & {\color{bkRed}\bfseries 10.0008890058…} & 10.000000000005 \\}
\bktabcaption{B列は、$h$ が小さすぎると10からはなれる（赤）。C列は10に近づいていく。}
\end{bkcenter}
$h=10^{-8}$ より小さくすると、B列の値は10からはなれてしまいます。コンピュータは数を16けたほどしか覚えていないので、とても近い2つの数をひき算すると、下のけたのずれが目立つのです（2-19）。$h$ の式 $10+5h$ を使えば、このずれは起きません。式で考えることの強みです。\par\vspace{4pt}
\textbf{ためしてみよう}\par
自分の表計算でも、$h$ を $\dfrac{1}{10}$ ずつ小さくして、B列が何回めから10からはなれはじめるか調べてみよう。
\end{column}

\clearpage
%% ---------------- 章末問題 ----------------
\bkendproblems
\bkproblemgroup{基本}{2}

\bkproblem 位置が $x=5t^2$（m）のボールについて、4秒から $4+\Delta t$ 秒までの平均の速さを、$\Delta t=0.1$ と $0.01$ で求めましょう。4秒の瞬間の速さは、いくつと考えられますか。
\bkpans{40.5（$x(4.1)=84.05$）、40.05（$x(4.01)=80.4005$）。\qty{40}{m/s}}

\bkproblem 次の式を展開しましょう。\quad（1）$(1+h)^2$\quad（2）$(3+h)^2$\quad（3）$(t+h)^2$
\bkpans{（1）\nobreak $1+2h+h^2$\quad（2）\nobreak $9+6h+h^2$\quad（3）\nobreak $t^2+2th+h^2$}

\bkproblem 位置が $x=4t^2$（m）の物体について、2秒から $2+h$ 秒までの平均の速さを $h$ の式で表し、2秒の瞬間の速さを求めましょう。
\bkpans{$4(2+h)^2=16+16h+4h^2$ より $16+4h$。\qty{16}{m/s}}

\bkproblem 位置が $x=t^2$（m）の物体の、5秒の瞬間の速さを求めましょう。
\bkpans{$(5+h)^2=25+10h+h^2$ より $10+h$。\qty{10}{m/s}}

\bkproblem 位置が $x=3t^2$（m）の物体の、$t$ 秒の瞬間の速さを $t$ の式で表しましょう。また、4秒の瞬間の速さを求めましょう。
\bkpans{$3(t+h)^2-3t^2=6th+3h^2$ より $v=6t$。4秒で\qty{24}{m/s}}

\bkproblem 位置が $x=2t^2+3t$（m）の物体の、1秒の瞬間の速さを求めましょう。
\bkpans{$x(1)=5$、$x(1+h)=5+7h+2h^2$ より $7+2h$。\qty{7}{m/s}}

\bkproblem 次の物体の、$t$ 秒の瞬間の速さを $v=2at$ を使って求めましょう。\quad（1）$x=6t^2$\quad（2）$x=0.5t^2$
\bkpans{（1）\nobreak $v=12t$\quad（2）\nobreak $v=t$}

\bkproblemgroup{チャレンジ}{2}

\bkproblem 位置が $x=3t^2+2t$（m）の物体の、2秒の瞬間の速さを求めましょう。
\bkpans{$x(2)=16$、$x(2+h)=16+14h+3h^2$ より $14+3h$。\qty{14}{m/s}}

\bkproblem $x=5t^2$（m）で落ちるボールの瞬間の速さは $v=10t$（m/s）でした。速さが\qty{25}{m/s}になるのは何秒後ですか。また、それまでに何m落ちていますか。
\bkpans{2.5秒後。$5\times2.5^2=\qty{31.25}{m}$}

\bkproblem 位置が $x=t^2-4t$（m）の物体があります。1秒の瞬間の速さを求めましょう。答えがマイナスになったら、それは何を表していますか。
\bkpans{$x(1)=-3$、$x(1+h)=-3-2h+h^2$ より $-2+h$。$\qty{-2}{m/s}$。$x$ が小さくなる向き（もどる向き）に、\qty{2}{m/s}で動いている}

\bkproblem 同じ所を同時に出発した2つの物体があります。Aの位置は $x=5t^2$、Bの位置は $x=20t$（m）です。2つの速さが同じになるのは何秒後ですか。そのとき、前にいるのはどちらで、何mはなれていますか。
\bkpans{Aの速さ $10t$、Bは\qty{20}{m/s}。$10t=20$ より2秒後。A $20$\,m、B $40$\,m で、Bが\qty{20}{m}前}

\begin{furikaeri}
\begin{checklist}
\item $\Delta t$ を小さくして、平均の速さを表で計算できる。\hfill{\footnotesize\sffamily\color{bkTheme}→\pageref{bk:5-2:hyou}ページ}
\item 平均の速さを $h$ の式で表せる。\hfill{\footnotesize\sffamily\color{bkTheme}→\pageref{bk:5-2:h}ページ}
\item $x=at^2$ の瞬間の速さを求められる。\hfill{\footnotesize\sffamily\color{bkTheme}→\pageref{bk:5-2:a}ページ}
\end{checklist}
\end{furikaeri}

\bkfukidashi{sora}{$h$ を0に「近づける」のと、0を「入れる」のはちがうんだね。次の章で、この考えに名前をつけるよ。}
